\documentclass[12pt,a4paper]{article}

\usepackage{fontawesome5}
\providecommand{\faPencilAlt}{\faIcon{pencil-alt}}

\usepackage{styles/couleurs}
\usepackage{styles/mise_en_page}
\usepackage{styles/boites}
\usepackage{styles/maths}
\usepackage{styles/figures}

\usepackage{enumitem}
\usepackage{array}
\usepackage{tabularx}
\usepackage{amssymb}
\usepackage{pdflscape}
\usepackage{tikz}
\usepackage{xparse}

\setlist[itemize]{leftmargin=1.4em,itemsep=0.15em,topsep=0.2em}
\setlist[enumerate]{leftmargin=*,itemsep=0.2em,topsep=0.2em}

\newcommand{\lignepointillee}{\noindent\dotfill\par\vspace{0.48em}}

\newcommand{\zonelignes}[1]{%
  \begin{tcolorbox}[
    enhanced,
    colback=white,
    colframe=NeutreC!40,
    boxrule=0.5pt,
    arc=1.5mm,
    left=3mm,
    right=3mm,
    top=2mm,
    bottom=1mm,
    before skip=0pt,
    after skip=0pt
  ]
    \foreach \i in {1,...,#1}{\lignepointillee}
  \end{tcolorbox}
}

\newcounter{exohc}

\NewDocumentCommand{\exohc}{m +m}{%
  \refstepcounter{exohc}
  \par\medskip
  \noindent
  \begin{minipage}[t]{0.47\linewidth}
    \vspace{0pt}
    {\sffamily\bfseries\color{TitreC}$\square$ Exercice \theexohc}
    \vspace{0.2em}
    #2
  \end{minipage}
  \hfill
  \begin{minipage}[t]{0.50\linewidth}
    \vspace{0pt}
    \zonelignes{#1}
  \end{minipage}
  \par\medskip
}

\begin{document}
\begin{landscape}

\begin{center}
  {\sffamily\LARGE\bfseries\color{TitreC}Exercices hors classe}\\[0.2em]
  {\sffamily\large\bfseries\color{ActionC}Arithmétique}
\end{center}

\remarque{Cette fiche est donnée à l'avance.
Ne faire que les exercices indiqués par le professeur.
Le mot d'ordre est : exercices courts, réguliers et correctement justifiés.}

\section*{A. Multiples, diviseurs et division euclidienne}

\exohc{5}{
Complète avec « multiple » ou « diviseur ».

\begin{enumerate}[label=\alph*)]
  \item \(8\) est un \ldots{} de \(56\).
  \item \(72\) est un \ldots{} de \(9\).
  \item \(5\) est un \ldots{} de \(45\).
  \item \(63\) est un \ldots{} de \(7\).
\end{enumerate}
}

\exohc{5}{
Écris les égalités de division euclidienne.

\[
173\div12
\qquad
458\div25
\qquad
1\,002\div17
\]
}

\exohc{5}{
Sans poser de division, indique si \(8\,370\) est divisible par :

\[
2,\qquad 3,\qquad 5,\qquad 9,\qquad 10.
\]

Justifie avec les critères de divisibilité.
}

\exohc{5}{
\begin{enumerate}[label=\alph*)]
  \item Donne les huit premiers multiples de \(6\).
  \item Donne tous les diviseurs de \(36\).
\end{enumerate}
}

\section*{B. Nombres premiers}

\exohc{5}{
Indique si chaque nombre est premier.
S'il ne l'est pas, donne un diviseur différent de \(1\) et de lui-même.

\[
31\qquad 49\qquad 57\qquad 71\qquad 87.
\]
}

\exohc{5}{
Écris tous les nombres premiers compris entre \(40\) et \(70\).
}

\exohc{6}{
Réponds par vrai ou faux en justifiant.

\begin{enumerate}[label=\alph*)]
  \item Tous les nombres impairs sont premiers.
  \item Le seul nombre premier pair est \(2\).
  \item Aucun multiple de \(5\), supérieur à \(5\), n'est premier.
\end{enumerate}
}

\section*{C. Décomposition en facteurs premiers}

\exohc{6}{
Décompose en produit de facteurs premiers :

\[
60\qquad 84\qquad 150.
\]
}

\exohc{6}{
Décompose en produit de facteurs premiers :

\[
126\qquad 196\qquad 330.
\]
}

\exohc{5}{
Une élève écrit :

\[
180=2\times3\times30.
\]

Est-ce une décomposition en produit de facteurs premiers ?
Explique puis termine la décomposition.
}

\section*{D. Fractions égales et simplification}

\exohc{6}{
Simplifie à l'aide des facteurs premiers :

\[
\frac{42}{63}
\qquad
\frac{70}{105}
\qquad
\frac{84}{126}.
\]
}

\exohc{6}{
Parmi les fractions suivantes, regroupe celles qui sont égales :

\[
\frac{6}{15}
\qquad
\frac{10}{25}
\qquad
\frac{14}{35}
\qquad
\frac{15}{36}.
\]
}

\exohc{5}{
Complète avec deux fractions égales à \(\dfrac{8}{20}\).

\[
\frac{8}{20}
=
\frac{\ldots}{\ldots}
=
\frac{\ldots}{\ldots}.
\]
}

\section*{E. Petits problèmes}

\exohc{7}{
On veut répartir \(54\) objets rouges et \(72\) objets bleus
dans des groupes identiques, sans reste.

\begin{enumerate}[label=\alph*)]
  \item Peut-on réaliser \(6\) groupes ?
  \item Peut-on réaliser \(9\) groupes ?
  \item Donne tous les nombres de groupes possibles.
\end{enumerate}
}

\exohc{7}{
Un panneau rectangulaire mesure \(240\) cm sur \(360\) cm.
On veut le recouvrir avec des carreaux carrés identiques,
sans découpe et sans espace.

\begin{enumerate}[label=\alph*)]
  \item Peut-on utiliser des carreaux de \(15\) cm de côté ?
  \item Peut-on utiliser des carreaux de \(14\) cm de côté ?
  \item Donne trois autres longueurs possibles pour le côté d'un carreau.
\end{enumerate}
}

\exohc{7}{
Je suis un entier compris entre \(100\) et \(400\).
Je suis pair, divisible par \(11\), par \(3\) et par \(5\).

Qui suis-je ? Justifie sans tester tous les entiers.
}

\end{landscape}
\end{document}
